\section{Introduction}
\label{sec:introduction}

A mobile robot runs perception next to mapping, planning, control, and logging,
usually on a single system-on-chip. Whichever processor perception occupies is
not available to the rest of the stack while it is busy, and whatever it draws
comes out of the same power budget. Embedded platforms therefore ship
fixed-function neural accelerators next to the GPU: on NVIDIA Jetson Orin
modules, up to two Deep Learning Accelerator (DLA) cores, depending on the SKU. Their appeal to a robot
integrator is not peak throughput; it is that a detector placed on one of them
stops competing for the GPU and, being fixed-function, is expected to convert
each frame into less energy than a general-purpose engine does.

The obstacle is that such an accelerator does not execute an arbitrary graph. It
implements an execution contract: a set of supported operators plus constraints
on tensor rank, layout, static shapes, precision, kernel parameters, and
dataflow. Regions that violate the contract either fail compilation or, if the
toolchain permits fallback, migrate to the shared processor, which reintroduces
exactly the contention the offload was meant to remove. Hardware support is
consequently not an exporter setting but a property of the model graph.

This paper studies the resulting decision as a systems question with a
measurable price. We take a strong reference detector, YOLO11n
\cite{jocher2024yolo11}, make it admissible with the smallest substitution we
could find that removes every contract violation, and then quantify what the
move to a fixed-function accelerator costs at the engine boundary, at the
application boundary, and in accuracy, so that the energy and availability
benefit can be weighed against a known price rather than assumed. The retrofit,
\emph{YOLO11-DLA-n}, keeps the YOLO11 stem, all C3k2 stages, SPPF, the complete
path-aggregation neck, the head towers, the losses, and the assignment rule, and
replaces two things only: the position-sensitive attention block, whose matrix
multiplications the contract forbids, and the decoded output interface, whose
rank-three anchor decode drags otherwise admissible producers off the
accelerator.

Stating the price plainly is the point of the paper, so we state it here. On the
target module, the strict DLA engine of the retrofit runs at $27.515$~ms median
engine-plus-transfer latency at batch one and $640\times640$, against $3.740$~ms
for the unmodified detector's GPU engine: fixed-function placement is
$7.36\times$ slower at the engine boundary, not faster. The substitution also
costs $1.72$ points of COCO $\mathrm{mAP}_{50\text{--}95}$. Neither number is a
result to be optimized away in this paper; together they are the price of
admission, and our claim is that in a power-capped robot whose GPU is
oversubscribed the price can be repaid in energy per frame and in restored GPU
availability.

Two intermediate positions deserve to be closed off early. The first is partial
offloading. TensorRT can place what it can on DLA and run the rest on the GPU;
for the unmodified detector this produces four DLA loadables, eleven GPU compute
kernels, and 22 reformat nodes, spends $94.0\%$ of engine time on DLA, and is
still slower than the fully placed retrofit and far slower than the GPU engine.
Fallback builds are useful as diagnostics and useless as deployments, so we
report the fallback build only in that role. The second is the temptation to
present accuracy retention as the contribution. Retaining $95.6\%$ of the
reference $\mathrm{mAP}_{50\text{--}95}$ is what makes the offload
\emph{admissible as a perception module at all}; it is a precondition, not a
finding, and we treat it as such.

The contributions are:
\begin{itemize}
  \item a formulation of fixed-function perception offloading as an
  admissibility problem over a declared hardware execution contract, with the
  learned graph separated from variable-cardinality task logic at an explicit
  accelerator boundary;
  \item a minimal admissible retrofit of a current one-stage detector under the
  NVIDIA DLA contract, with a weight-aware initialization that transfers the
  reference model's projections into the substituted block;
  \item a measured price of admission on the target module: strict-build
  acceptance, engine-boundary latency and throughput, accuracy retention, and a
  stage-resolved application-level accounting that separates the accelerator's
  cost from the reference implementation's cost; and
  \item evidence that partial offloading does not provide a compromise, plus the
  protocol and instrumentation for the energy and GPU-isolation measurements on
  which the decision ultimately rests.
\end{itemize}

Hereafter DLA denotes NVIDIA's accelerator, not the homonymous Deep Layer
Aggregation architecture.
